\section{RSP：从 Historical ASL 到 Capability-Triggered Safeguards}

课堂使用 AI Safety Levels 描述模型能力与防护等级的关联，并强调 defense in depth。当前 Anthropic RSP v3.0 已演化为 capability thresholds、required safeguards、risk reports 与治理流程。两版共同的 durable idea 是：随着 capability evidence 接近高影响 threshold，训练和部署必须先升级 security/deployment controls，而不是事后补救。

\subsection{Threshold、safeguard 与 gate}

团队先通过 evaluation、elicitation 与 forecast 形成 capability evidence，再判断离 threshold 多远；若达到或接近 threshold，就要满足对应 required safeguards，例如更强 model-weight security、access control、runtime guard、monitoring、red team、incident readiness 与 governance review；最后才决定继续训练、扩大部署或限制 access。

\lecturefigure{12-capability-safeguards.png}{能力证据触发 threshold review、required safeguards 与 deployment gate。}{课堂 ASL framing；Anthropic RSP v3.0；概念重绘。}

\begin{knowledgebox}{读图：版本变化说明 governance 需要可修订}
课堂 ASL 把模型/防护分级绑定得较紧；RSP 后续版本把 capability threshold 与 safeguard package 拆开，以便针对具体风险更新。这个变化不是“旧版错误、新版正确”的简单替换，而是承认 evaluation science、threat model 与 safeguard evidence 会变化，因此 policy 本身也必须 versioned、公开变更并接受 review。
\end{knowledgebox}

\teachervoice{讲者强调评估最难之处之一是 elicitation overhang：今天没测出来，不代表更强 scaffold、tool 或研究 effort 后仍做不到。课堂里的 ASL 不是给模型贴永久标签，而是要求组织在不确定性下提前准备防护。}

\subsection{Defense in Depth}

\term{defense in depth}（纵深防御）是设置相互独立、失败模式不同的多层控制。Training 通过 data/objective 与 post-training 降低危险行为；access 层限制谁能使用什么能力；runtime classifier/tool policy 监控请求和 action；response 层负责 red team、incident、patch、revoke 与 customer communication。

\lecturefigure{13-defense-in-depth.png}{Frontier model safety 需要 training、access、runtime 与 response 的纵深防御。}{本地字幕 37:00--39:10；概念重绘。}

\begin{knowledgebox}{读图：层数多不等于独立}
若 training、classifier 和 judge 都依赖同一模型与同一数据，它们可能共享 blind spot。有效纵深防御要分析 common-mode failure、权限分离、监控独立性和 human override。每层都应产生 evidence 供下一层判断，并有明确 fallback，而不是把风险在层间循环转发。
\end{knowledgebox}

\subsection{本章小结}

RSP 的核心是 capability-triggered governance。术语会变，durable contract 是：先评估，再满足 required safeguards，最后训练/部署；任何例外都要有 owner、证据与重新评估日期。

